\begin{statement}[C.8.1; Proved: three uniform bounds on the elliptic norm]
For every $z\in\C$ and every $\tau$ in the standard modular
fundamental region, we have
\[
 \begin{array}{c|ccc}
 \text{restriction}&T\le2&T\le7/5&T\le11/10\\ \hline
 N(z,\tau)&<17/10&<3/2&<143/100.
 \end{array}
\]
If $\wp_\tau(z)=0$ and $T\ge2$, then $N(z,\tau)<541/500$.
\end{statement}
\begin{proof}
Taking absolute values in the theta series and completing the square
gives
\[
 e^{-\pi(\Im z)^2/T}|\vartheta_1(z,\tau)|
 \le\sum_{n\in\Z}
 e^{-\pi T(n+1/2+\Im z/T)^2}
 \le\sum_{n\in\Z}e^{-\pi Tn^2}.
\]
The last inequality follows from Poisson summation: the Fourier
coefficients of the periodized Gaussian are positive, so its maximum
is at an integer. The eta product and $n^2\ge1+3(n-1)$ for $n\ge1$
therefore imply, on $a\le T\le b$,
\begin{equation}\label{eq:ca-sup}
 N(z,\tau)\le
 \exp\!\left(\frac{r_a}{(1-r_a)^2}\right)
 \max\{F(a),F(b)\},\qquad
 F(T)=e^{\pi T/12}
 \left(1+\frac{2e^{-\pi T}}{1-e^{-3\pi T}}\right),
\end{equation}
where $r_a=e^{-2\pi a}$. Indeed,
$-\sum\log(1-r^n)\le r/(1-r)^2$. The function $F$ is convex,
as follows by expanding its denominator as a geometric series;
each summand is an exponential of a real linear function of $T$.
We use $6/7<\sqrt3/2$ and the following rational bounds:
\[
\begin{array}{c|@{\,}c@{\,}c@{\,}c@{\,}c@{\,}c}
T&6/7&1&11/10&7/5&2\\ \hline
F(T)&<14211/10000&<14116/10000&<14180/10000&<14782/10000&<16944/10000.
\end{array}
\]
The eta factors in \eqref{eq:ca-sup} are less than $10047/10000$
for $a=6/7$ and less than $10019/10000$ for $a=1$.
Multiplying these fractions proves the asserted bounds, using
$[6/7,1]$ and $[1,2]$ separately for the first one.

We specify elementary arithmetic for all the scalar estimates in
this proof. One may use $314159/100000<\pi<355/113$ and, for
$0\le x\le20$, the Taylor bounds
\begin{equation}\label{eq:ca-exp}
 S_{80}(x)\le e^x\le
 S_{80}(x)+\frac{x^{81}}{81!}\frac1{1-x/82},
 \qquad S_{80}(x)=\sum_{j=0}^{80}\frac{x^j}{j!}.
\end{equation}
Taking reciprocals bounds negative exponentials. The remainder
bound follows by comparison of consecutive terms, not by numerical
sampling. These inequalities verify each displayed rational bound;
much shorter Taylor sums also suffice. Thus the decimal diagnostics
in the accompanying programs are dispensable.

For the cusp we check directly the analytic step of supplied K6.
Choose $|\Im z|\le T/2$ by translation, put $b_0=43/100$ and
$r_0=e^{-4\pi}$. The Fourier series
\[
 \frac{\wp_\tau(z)}{\pi^2}
 =\csc^2(\pi z)-\frac13
   +8\sum_{n\ge1}\frac{nq^n}{1-q^n}
          \bigl(1-\cos(2\pi nz)\bigr)
\]
shows, if $|\Im z|\ge b_0$, that
\begin{align*}
 \left|\frac{\wp_\tau(z)}{\pi^2}+\frac13\right|
 &\le \frac1{\sinh^2(\pi b_0)}
 +8\left\{\frac{r_0}{(1-r_0)^3}
 +\frac1{2(1-r_0)}
 \left(\frac{e^{-2\pi}}{(1-e^{-2\pi})^2}
       +\frac{e^{-6\pi}}{(1-e^{-6\pi})^2}\right)\right\}
 <\frac13.
\end{align*}
The two terms in the last parentheses bound respectively
$r^n e^{2\pi n|\Im z|}$ and $r^n e^{-2\pi n|\Im z|}$.
This inequality excludes a zero. Hence a zero has
$|\Im z|<b_0$. The Jacobi product, $\log(1+x)\le x$, and
the monotonicity of $\cosh(\pi y)e^{-\pi y^2/T}$ on
$0\le y\le b_0$ give
\[
 N(z,\tau)\le
 2\cosh(\pi b_0)e^{-\pi/3-\pi b_0^2/2}
 \exp\!\left(
 \frac{2\cosh(2\pi b_0)r_0}{1-r_0}
 +\frac{r_0^2}{1-r_0^2}\right)<\frac{541}{500}.
\]
For the stated monotonicity, the concave function
$\tanh(\pi y)-y$ is zero at zero and positive at $b_0$.
For $T\ge2$ the right side before substituting $T=2$ decreases:
the derivative of $-\pi T/6-\pi b_0^2/T$ is negative and
both product errors decrease. The two strict scalar inequalities
follow from \eqref{eq:ca-exp}. This proves the cusp assertion
without using any certificate from K6.
\end{proof}
